% Source printed p. 414 / PDF p. 422.
\setlength{\emergencystretch}{4em}
\chapter{几何不变量理论}\label{src:8}
\hypersetup{pageanchor=false}
\AddToHook{shipout/foreground}[chapter-eight-index-pages]{\put(0,0){\special{pdf:dest (page.\thepage) [@thispage /Fit]}}}
\begin{epigraph}
有时很难体会某些定义和定理在六十年前曾带来多么深刻的变革. Mumford 的若干思想对他的前辈而言似乎十分陌生, 但他一经提出, 它们很快就被视为``显而易见的方法''.
\epigraphauthor{János Kollár \cite[第 623 页]{Kol21}}
\end{epigraph}
Mumford 在 \cite{GIT} 中发展了\newterm{几何不变量理论}{Geometric Invariant Theory} (GIT), 用以构造代数几何中的商和模空间. 关于其他阐述, 我们推荐 \cite{New78,Kra84,Dol03,Muk03,Stu08}.
